\section{Introduction and Positioning}
How does matched connection loss in a connectome-derived spiking neural network (SNN) affect first contact as feedback, observation horizon and calibrated sensing change? We pair healthy and damaged runs, resolving neural commands, action timing and contact separately.

Shiu's connectome-derived LIF model\cite{shiu2024drosophila} supplies commands to our synthetic pursuit--escape plant. An\cite{an2026sparse} studies task/state-dependent sparsification; Jo\cite{jo2026simulating} links damaged circuits to virtual-arena behavior. We quantify paired first contact under frozen laws and horizons, and test future-contact warnings from a fixed 400-ms prefix.

DriveFI searches closed-loop safety violations\cite{jha2019drivefi}; SpikeFI, SpikingJET and enpheeph provide SNN fault injection\cite{spyrou2024spikefi,gogebakan2024spikingjet,colucci2022enpheeph}. Common random numbers\cite{glasserman1992crn}, trajectory monitoring\cite{donze2010breach}, benchmarking\cite{yik2023neurobench,dangelo2026embodied} and physical fault propagation\cite{chessa2026crosslayer} are established. Our ablation tests which pairing, timing and native-replay obligations add findings beyond trusted record identities.

We contribute: (i) stage-resolved paired timing/contact measurements, including law-discordant outcomes under shared schedules and horizon-sensitive, uncertain RF transfer; (ii) frozen-prefix warnings with 11/13 missed RF-0 contacts and cross-configuration false alarms; and (iii) a frozen-checker ablation quantifying complementary findings, misses, costs and historical reporting bias across fly and tank adapters. Timing remains descriptive; the engineering comparison tests audit information, not neural authenticity.
